\documentclass[12pt,a4paper]{article}
\usepackage[utf8]{inputenc}
\usepackage[T1]{fontenc}
\usepackage{lmodern}
\usepackage[margin=1in]{geometry}
\usepackage{hyperref}
\usepackage{xcolor}
\usepackage{listings}
\usepackage{tcolorbox}
\usepackage{enumitem}
\usepackage{fancyhdr}
\usepackage{titlesec}
\usepackage{graphicx}
\usepackage{amsmath}
\usepackage{amssymb}
\usepackage{amsfonts}
\usepackage{fontawesome5}
\usepackage{tikz}
\usetikzlibrary{positioning,arrows.meta,shapes}

% Color definitions
\definecolor{codegreen}{rgb}{0,0.6,0}
\definecolor{codegray}{rgb}{0.5,0.5,0.5}
\definecolor{codepurple}{rgb}{0.58,0,0.82}
\definecolor{backcolour}{rgb}{0.97,0.97,0.97}
\definecolor{checkpoint}{rgb}{0.2,0.6,0.2}
\definecolor{gitcolor}{rgb}{0.9,0.4,0.1}
\definecolor{warningcolor}{rgb}{0.8,0.2,0.2}
\definecolor{conceptcolor}{rgb}{0.1,0.3,0.6}
\definecolor{mathcolor}{rgb}{0.5,0.0,0.5}

% Code listing style
\lstdefinestyle{mystyle}{
    backgroundcolor=\color{backcolour},
    commentstyle=\color{codegreen},
    keywordstyle=\color{blue}\bfseries,
    numberstyle=\tiny\color{codegray},
    stringstyle=\color{codepurple},
    basicstyle=\ttfamily\small,
    breakatwhitespace=false,
    breaklines=true,
    captionpos=b,
    keepspaces=true,
    numbers=left,
    numbersep=5pt,
    showspaces=false,
    showstringspaces=false,
    showtabs=false,
    tabsize=2,
    frame=single,
    rulecolor=\color{codegray}
}
\lstset{style=mystyle}

% Custom environments
\newtcolorbox{checkpoint}{
    colback=green!5!white,
    colframe=checkpoint,
    title={\faCheckCircle\ \ CHECKPOINT: Verify Your Work},
    fonttitle=\bfseries
}

\newtcolorbox{gitcommit}{
    colback=orange!5!white,
    colframe=gitcolor,
    title={\faGitAlt\ \ GIT COMMIT REQUIRED},
    fonttitle=\bfseries
}

\newtcolorbox{concept}{
    colback=blue!5!white,
    colframe=conceptcolor,
    title={\faLightbulb\ \ FUNDAMENTAL CONCEPT},
    fonttitle=\bfseries
}

\newtcolorbox{warning}{
    colback=red!5!white,
    colframe=warningcolor,
    title={\faExclamationTriangle\ \ COMMON MISTAKE},
    fonttitle=\bfseries
}

\newtcolorbox{mathbox}{
    colback=purple!5!white,
    colframe=mathcolor,
    title={\faSquareRootAlt\ \ MATHEMATICAL FOUNDATION},
    fonttitle=\bfseries
}

\newtcolorbox{explanation}{
    colback=gray!5!white,
    colframe=codegray,
    title={\faCode\ \ LINE-BY-LINE EXPLANATION},
    fonttitle=\bfseries
}

% Header/Footer
\pagestyle{fancy}
\fancyhf{}
\rhead{Graph Embeddings for Retrieval}
\lhead{First Principles Tutorial}
\cfoot{\thepage}

\title{\Huge\textbf{Building a Retrieval System with Graph Embeddings}\\[0.5cm]
\Large A First-Principles Guide from Linear Algebra to Production\\[0.3cm]
\normalsize From Zero to Semantic Search}
\author{Comprehensive Tutorial for Machine Learning Engineers}
\date{January 2026}

\begin{document}

\maketitle
\thispagestyle{empty}

\vspace{1cm}

\begin{abstract}
This document provides an exhaustive, line-by-line guide to building a retrieval system using graph embeddings. Starting from the mathematics of vectors and graphs, we progress through neural network fundamentals, graph neural networks, and finally implement a complete semantic retrieval system. Every equation is derived, every line of code is explained, and every concept is built from first principles. No prior machine learning knowledge is assumed---only basic Python familiarity.
\end{abstract}

\newpage
\tableofcontents
\newpage

%==============================================================================
% PART I: MATHEMATICAL FOUNDATIONS
%==============================================================================
\part{Mathematical Foundations}

\section{Chapter 1: What is Retrieval?}

Before writing any code, we must understand what we're building.

\begin{concept}
\textbf{Retrieval} is the task of finding relevant items from a large collection given a query. Examples:
\begin{itemize}
    \item Google Search: query = ``best pizza NYC'', retrieve = relevant web pages
    \item Netflix: query = your watch history, retrieve = movies you'll like
    \item Research: query = ``graph neural networks'', retrieve = relevant papers
\end{itemize}
A retrieval system answers: ``Given this query, what items are most relevant?''
\end{concept}

\subsection{The Evolution of Retrieval}

\subsubsection{Stage 1: Keyword Matching (1970s-1990s)}

The simplest approach: if the query word appears in a document, it's relevant.

\begin{lstlisting}[language=Python,caption={Naive keyword matching}]
def keyword_search(query, documents):
    results = []
    for doc in documents:
        if query.lower() in doc.lower():
            results.append(doc)
    return results
\end{lstlisting}

\textbf{Problem}: ``dog food'' won't match ``canine nutrition'' even though they're semantically related.

\subsubsection{Stage 2: TF-IDF (1970s-2000s)}

\begin{concept}
\textbf{TF-IDF} (Term Frequency - Inverse Document Frequency) weighs words by:
\begin{itemize}
    \item \textbf{TF}: How often the word appears in this document (more = more relevant)
    \item \textbf{IDF}: How rare the word is across all documents (rarer = more distinctive)
\end{itemize}
\end{concept}

\begin{mathbox}
\textbf{TF-IDF Formula:}
\[
\text{TF-IDF}(t, d, D) = \text{TF}(t, d) \times \text{IDF}(t, D)
\]
Where:
\[
\text{TF}(t, d) = \frac{\text{count of term } t \text{ in document } d}{\text{total terms in } d}
\]
\[
\text{IDF}(t, D) = \log\left(\frac{\text{total documents in } D}{\text{documents containing } t}\right)
\]
\end{mathbox}

\textbf{Problem}: Still relies on exact word matching. ``dog'' and ``puppy'' are treated as completely different.

\subsubsection{Stage 3: Semantic Embeddings (2013-Present)}

\begin{concept}
\textbf{Embeddings} represent words, sentences, or documents as vectors (lists of numbers) in a continuous space where \textit{similar meanings are close together}.

Instead of matching words, we match \textit{meaning}.
\end{concept}

This tutorial focuses on the most powerful form: \textbf{Graph Embeddings}---where we leverage relationships between items to create even better representations.

%==============================================================================
% CHAPTER 2: VECTORS AND LINEAR ALGEBRA
%==============================================================================
\section{Chapter 2: Vectors and Linear Algebra}

To understand embeddings, you must understand vectors.

\subsection{What is a Vector?}

\begin{concept}
A \textbf{vector} is an ordered list of numbers. It can represent:
\begin{itemize}
    \item A point in space: $[3, 4]$ is a point at x=3, y=4
    \item A direction: from origin to that point
    \item Features: $[\text{height}, \text{weight}, \text{age}]$ describes a person
\end{itemize}
\end{concept}

\begin{lstlisting}[language=Python,caption={Vectors in Python}]
import numpy as np

# A 2D vector (point in 2D space)
v1 = np.array([3, 4])

# A 3D vector (point in 3D space)
v2 = np.array([1, 2, 3])

# A 768-dimensional vector (like BERT embeddings)
v3 = np.random.randn(768)

print(f"v1 shape: {v1.shape}")  # (2,)
print(f"v2 shape: {v2.shape}")  # (3,)
print(f"v3 shape: {v3.shape}")  # (768,)
\end{lstlisting}

\begin{explanation}
\textbf{Line 1:} Import NumPy, the fundamental library for numerical computing in Python.

\textbf{Line 4:} Create a 2D vector using \texttt{np.array()}. This is a point at coordinates (3, 4).

\textbf{Line 7:} A 3D vector---like a point in physical space (x, y, z).

\textbf{Line 10:} A 768-dimensional vector. We can't visualize this, but mathematically it works the same way. \texttt{np.random.randn(768)} creates 768 random numbers from a standard normal distribution.

\textbf{Line 12:} \texttt{.shape} tells us the dimensions. \texttt{(2,)} means a 1D array with 2 elements.
\end{explanation}

\subsection{Vector Operations}

\subsubsection{Vector Addition}

\begin{mathbox}
For vectors $\mathbf{a} = [a_1, a_2, ..., a_n]$ and $\mathbf{b} = [b_1, b_2, ..., b_n]$:
\[
\mathbf{a} + \mathbf{b} = [a_1 + b_1, a_2 + b_2, ..., a_n + b_n]
\]
\end{mathbox}

\begin{lstlisting}[language=Python,caption={Vector addition}]
a = np.array([1, 2, 3])
b = np.array([4, 5, 6])

c = a + b
print(c)  # [5, 7, 9]
\end{lstlisting}

\subsubsection{Scalar Multiplication}

\begin{mathbox}
Multiplying a vector by a scalar (single number):
\[
k \cdot \mathbf{a} = [k \cdot a_1, k \cdot a_2, ..., k \cdot a_n]
\]
\end{mathbox}

\begin{lstlisting}[language=Python,caption={Scalar multiplication}]
a = np.array([1, 2, 3])
scaled = 2 * a
print(scaled)  # [2, 4, 6]
\end{lstlisting}

\subsubsection{Vector Magnitude (Length)}

\begin{mathbox}
The magnitude (or norm) of a vector is its ``length'':
\[
||\mathbf{a}|| = \sqrt{a_1^2 + a_2^2 + ... + a_n^2}
\]
This is the Euclidean norm, also called L2 norm.
\end{mathbox}

\begin{lstlisting}[language=Python,caption={Vector magnitude}]
a = np.array([3, 4])

# Manual calculation
magnitude_manual = np.sqrt(3**2 + 4**2)
print(magnitude_manual)  # 5.0

# Using numpy
magnitude_np = np.linalg.norm(a)
print(magnitude_np)  # 5.0
\end{lstlisting}

\begin{explanation}
\textbf{Line 4:} Manual calculation: $\sqrt{3^2 + 4^2} = \sqrt{9 + 16} = \sqrt{25} = 5$

\textbf{Line 8:} \texttt{np.linalg.norm()} computes the Euclidean norm automatically.
\end{explanation}

\subsection{The Dot Product---The Heart of Similarity}

\begin{mathbox}
The \textbf{dot product} of two vectors:
\[
\mathbf{a} \cdot \mathbf{b} = a_1 b_1 + a_2 b_2 + ... + a_n b_n = \sum_{i=1}^{n} a_i b_i
\]

Geometrically:
\[
\mathbf{a} \cdot \mathbf{b} = ||\mathbf{a}|| \cdot ||\mathbf{b}|| \cdot \cos(\theta)
\]
where $\theta$ is the angle between the vectors.
\end{mathbox}

\begin{concept}
The dot product measures \textbf{how aligned} two vectors are:
\begin{itemize}
    \item Positive dot product: vectors point in similar directions
    \item Zero dot product: vectors are perpendicular (orthogonal)
    \item Negative dot product: vectors point in opposite directions
\end{itemize}
\end{concept}

\begin{lstlisting}[language=Python,caption={Dot product}]
a = np.array([1, 0])  # Points right
b = np.array([1, 0])  # Points right (same direction)
c = np.array([0, 1])  # Points up (perpendicular)
d = np.array([-1, 0]) # Points left (opposite)

print(np.dot(a, b))  # 1 (same direction)
print(np.dot(a, c))  # 0 (perpendicular)
print(np.dot(a, d))  # -1 (opposite direction)
\end{lstlisting}

\subsection{Cosine Similarity---Normalized Comparison}

\begin{mathbox}
\textbf{Cosine similarity} measures the angle between vectors, ignoring their magnitudes:
\[
\text{cos\_sim}(\mathbf{a}, \mathbf{b}) = \frac{\mathbf{a} \cdot \mathbf{b}}{||\mathbf{a}|| \cdot ||\mathbf{b}||} = \cos(\theta)
\]

Range: $[-1, 1]$ where:
\begin{itemize}
    \item $1$ = identical direction (most similar)
    \item $0$ = perpendicular (unrelated)
    \item $-1$ = opposite direction (most dissimilar)
\end{itemize}
\end{mathbox}

\begin{lstlisting}[language=Python,caption={Cosine similarity implementation}]
def cosine_similarity(a, b):
    """
    Compute cosine similarity between two vectors.

    Args:
        a: First vector (numpy array)
        b: Second vector (numpy array)

    Returns:
        Cosine similarity score between -1 and 1
    """
    dot_product = np.dot(a, b)
    norm_a = np.linalg.norm(a)
    norm_b = np.linalg.norm(b)

    # Avoid division by zero
    if norm_a == 0 or norm_b == 0:
        return 0.0

    return dot_product / (norm_a * norm_b)

# Test it
v1 = np.array([1, 2, 3])
v2 = np.array([1, 2, 3])  # Identical
v3 = np.array([3, 2, 1])  # Different

print(cosine_similarity(v1, v2))  # 1.0 (identical)
print(cosine_similarity(v1, v3))  # ~0.71 (somewhat similar)
\end{lstlisting}

\begin{explanation}
\textbf{Lines 12-14:} Calculate the three components: dot product and both norms.

\textbf{Lines 17-18:} Edge case handling---if either vector has zero magnitude, we can't divide by it.

\textbf{Line 20:} The actual formula: dot product divided by product of norms.
\end{explanation}

\begin{concept}
\textbf{Why Cosine Similarity for Retrieval?}

Cosine similarity is the standard metric for comparing embeddings because:
\begin{enumerate}
    \item It's \textbf{scale-invariant}: $[1, 2]$ and $[100, 200]$ have similarity 1.0
    \item It captures \textbf{direction}, which represents meaning
    \item It's \textbf{fast to compute} with optimized libraries
\end{enumerate}
\end{concept}

\begin{checkpoint}
Open a Python interpreter and verify:
\begin{lstlisting}[language=Python]
import numpy as np
a = np.array([1, 0, 0])
b = np.array([0, 1, 0])
print(np.dot(a, b))  # Should print 0
\end{lstlisting}
If you get 0, you understand that perpendicular vectors have zero dot product.
\end{checkpoint}

%==============================================================================
% CHAPTER 3: WHAT ARE GRAPHS?
%==============================================================================
\section{Chapter 3: Introduction to Graphs}

\begin{concept}
A \textbf{graph} is a mathematical structure consisting of:
\begin{itemize}
    \item \textbf{Nodes} (vertices): The entities/objects
    \item \textbf{Edges}: Connections/relationships between nodes
\end{itemize}

Graphs model relationships: social networks (people connected by friendships), citation networks (papers connected by references), knowledge graphs (entities connected by relations).
\end{concept}

\subsection{Graph Notation}

\begin{mathbox}
A graph $G = (V, E)$ where:
\begin{itemize}
    \item $V = \{v_1, v_2, ..., v_n\}$ is the set of $n$ nodes
    \item $E \subseteq V \times V$ is the set of edges
\end{itemize}

An edge $(v_i, v_j) \in E$ means node $v_i$ is connected to node $v_j$.
\end{mathbox}

\subsection{Types of Graphs}

\begin{enumerate}
    \item \textbf{Undirected}: Edges have no direction. If A connects to B, then B connects to A.
    \begin{itemize}
        \item Example: Facebook friendships (mutual)
    \end{itemize}

    \item \textbf{Directed}: Edges have direction. A$\rightarrow$B doesn't imply B$\rightarrow$A.
    \begin{itemize}
        \item Example: Twitter follows (you can follow someone who doesn't follow you back)
    \end{itemize}

    \item \textbf{Weighted}: Edges have numerical weights representing strength.
    \begin{itemize}
        \item Example: Road network (edge weight = distance)
    \end{itemize}

    \item \textbf{Attributed}: Nodes and/or edges have feature vectors.
    \begin{itemize}
        \item Example: Social network where each user node has features like [age, location, interests]
    \end{itemize}
\end{enumerate}

\subsection{Graph Representations}

\subsubsection{Adjacency Matrix}

\begin{mathbox}
The \textbf{adjacency matrix} $A$ is an $n \times n$ matrix where:
\[
A_{ij} = \begin{cases}
1 & \text{if edge exists from } v_i \text{ to } v_j \\
0 & \text{otherwise}
\end{cases}
\]

For undirected graphs, $A$ is symmetric: $A_{ij} = A_{ji}$
\end{mathbox}

\begin{lstlisting}[language=Python,caption={Adjacency matrix representation}]
import numpy as np

# A simple graph with 4 nodes
# Edges: 0-1, 0-2, 1-2, 2-3
adj_matrix = np.array([
    [0, 1, 1, 0],  # Node 0 connects to 1 and 2
    [1, 0, 1, 0],  # Node 1 connects to 0 and 2
    [1, 1, 0, 1],  # Node 2 connects to 0, 1, and 3
    [0, 0, 1, 0],  # Node 3 connects to 2
])

print("Adjacency Matrix:")
print(adj_matrix)

# Check if edge exists between node 0 and node 1
print(f"\nEdge 0-1 exists: {adj_matrix[0, 1] == 1}")  # True
print(f"Edge 0-3 exists: {adj_matrix[0, 3] == 1}")  # False
\end{lstlisting}

\begin{explanation}
\textbf{Lines 5-10:} Create a 4$\times$4 matrix. Row $i$, column $j$ being 1 means there's an edge from node $i$ to node $j$.

\textbf{Line 16:} Access matrix element at position [0, 1]---checking if nodes 0 and 1 are connected.
\end{explanation}

\subsubsection{Edge List}

\begin{lstlisting}[language=Python,caption={Edge list representation}]
# Same graph as edge list
edges = [
    (0, 1),
    (0, 2),
    (1, 2),
    (2, 3),
]

# For undirected graphs, we often store both directions
edges_undirected = [
    (0, 1), (1, 0),
    (0, 2), (2, 0),
    (1, 2), (2, 1),
    (2, 3), (3, 2),
]
\end{lstlisting}

\subsubsection{Adjacency List}

\begin{lstlisting}[language=Python,caption={Adjacency list representation}]
# Adjacency list: dictionary mapping node to its neighbors
adj_list = {
    0: [1, 2],
    1: [0, 2],
    2: [0, 1, 3],
    3: [2],
}

# Get neighbors of node 2
print(f"Neighbors of node 2: {adj_list[2]}")  # [0, 1, 3]
\end{lstlisting}

\begin{concept}
\textbf{When to Use Each Representation:}
\begin{itemize}
    \item \textbf{Adjacency Matrix}: Dense graphs, when you need fast edge lookup $O(1)$, matrix operations
    \item \textbf{Edge List}: Sparse graphs, easy to iterate over all edges
    \item \textbf{Adjacency List}: Sparse graphs, fast neighbor lookup
\end{itemize}
\end{concept}

\subsection{Graph Metrics}

\begin{lstlisting}[language=Python,caption={Basic graph metrics}]
def degree(adj_matrix, node):
    """Number of edges connected to a node."""
    return np.sum(adj_matrix[node])

def neighbors(adj_list, node):
    """Get all neighbors of a node."""
    return adj_list.get(node, [])

def num_edges(adj_matrix):
    """Total number of edges (undirected graph)."""
    return np.sum(adj_matrix) // 2  # Divide by 2 for undirected

# Examples
print(f"Degree of node 2: {degree(adj_matrix, 2)}")  # 3
print(f"Total edges: {num_edges(adj_matrix)}")  # 4
\end{lstlisting}

\begin{mathbox}
\textbf{Degree} of a node = number of edges connected to it.
\[
\text{deg}(v_i) = \sum_{j=1}^{n} A_{ij}
\]
\end{mathbox}

%==============================================================================
% CHAPTER 4: WHY GRAPH EMBEDDINGS?
%==============================================================================
\section{Chapter 4: Why Graph Embeddings?}

\begin{concept}
\textbf{The Core Problem:}

Graphs are \textit{discrete} structures. Machine learning algorithms (neural networks, regression, clustering) work with \textit{continuous} vectors.

\textbf{Graph embeddings} convert graph structure into continuous vectors while preserving important properties like:
\begin{itemize}
    \item Which nodes are similar
    \item Which nodes are connected
    \item The overall structure/community
\end{itemize}
\end{concept}

\subsection{What Should Good Embeddings Capture?}

\begin{enumerate}
    \item \textbf{Structural similarity}: Nodes with similar connections should have similar embeddings
    \item \textbf{Proximity}: Connected nodes should be close in embedding space
    \item \textbf{Community structure}: Nodes in the same cluster should be close
    \item \textbf{Node features}: If nodes have attributes, incorporate them
\end{enumerate}

\subsection{The Embedding Space}

\begin{mathbox}
An \textbf{embedding function} maps nodes to vectors:
\[
f: V \rightarrow \mathbb{R}^d
\]
Where $d$ is the embedding dimension (typically 64, 128, 256, or 768).

For each node $v_i$, we get an embedding vector $\mathbf{z}_i \in \mathbb{R}^d$.
\end{mathbox}

\begin{lstlisting}[language=Python,caption={Conceptual embedding lookup}]
# Imagine we have learned embeddings for 1000 nodes
# Each embedding is a 128-dimensional vector
num_nodes = 1000
embedding_dim = 128

# Embedding matrix: each row is a node's embedding
embeddings = np.random.randn(num_nodes, embedding_dim)

# Get embedding for node 42
node_42_embedding = embeddings[42]
print(f"Node 42 embedding shape: {node_42_embedding.shape}")  # (128,)

# Find similarity between node 42 and node 100
similarity = cosine_similarity(embeddings[42], embeddings[100])
print(f"Similarity between node 42 and 100: {similarity:.4f}")
\end{lstlisting}

\subsection{Graph Embeddings for Retrieval}

\begin{concept}
\textbf{The Big Idea for Retrieval:}

\begin{enumerate}
    \item Represent your items (documents, products, users) as nodes in a graph
    \item Connect nodes based on relationships (co-purchases, citations, hyperlinks)
    \item Learn embeddings that capture both node features AND graph structure
    \item For retrieval: find nodes whose embeddings are closest to the query embedding
\end{enumerate}

This is more powerful than simple text embeddings because you leverage \textit{relational information}!
\end{concept}

\begin{checkpoint}
Create a simple graph in Python:
\begin{lstlisting}[language=Python]
adj = {0: [1, 2], 1: [0, 2], 2: [0, 1, 3], 3: [2]}
print(f"Node 2 has {len(adj[2])} neighbors")  # Should print 3
\end{lstlisting}
\end{checkpoint}

%==============================================================================
% PART II: EMBEDDING METHODS
%==============================================================================
\part{Graph Embedding Methods}

\section{Chapter 5: Classical Methods - Random Walk Embeddings}

\subsection{The Random Walk Intuition}

\begin{concept}
A \textbf{random walk} is a sequence of nodes where each step moves to a random neighbor.

The key insight: Nodes that appear in similar random walks are structurally similar.

If random walks from node A often visit nodes X, Y, Z, and random walks from node B also often visit X, Y, Z, then A and B are similar.
\end{concept}

\subsection{DeepWalk (2014)}

\begin{mathbox}
\textbf{DeepWalk Algorithm:}
\begin{enumerate}
    \item For each node $v$, generate multiple random walks starting from $v$
    \item Treat each walk as a ``sentence'' (sequence of node IDs)
    \item Apply Word2Vec (Skip-gram) to learn node embeddings
\end{enumerate}

The objective: Maximize probability of observing a node's neighbors given its embedding:
\[
\max_f \sum_{v \in V} \sum_{u \in N(v)} \log P(u | f(v))
\]
\end{mathbox}

\begin{lstlisting}[language=Python,caption={Random walk generation}]
import random

def random_walk(adj_list, start_node, walk_length):
    """
    Generate a random walk starting from start_node.

    Args:
        adj_list: Dictionary mapping node to list of neighbors
        start_node: Node to start the walk from
        walk_length: Number of steps in the walk

    Returns:
        List of nodes visited in the walk
    """
    walk = [start_node]
    current = start_node

    for _ in range(walk_length - 1):
        neighbors = adj_list.get(current, [])
        if not neighbors:
            break  # Dead end
        current = random.choice(neighbors)
        walk.append(current)

    return walk

# Example
adj_list = {0: [1, 2], 1: [0, 2], 2: [0, 1, 3], 3: [2]}

# Generate 3 random walks from node 0
for i in range(3):
    walk = random_walk(adj_list, start_node=0, walk_length=5)
    print(f"Walk {i+1}: {walk}")
\end{lstlisting}

\begin{explanation}
\textbf{Line 14:} Initialize the walk with the starting node.

\textbf{Line 17:} Loop for \texttt{walk\_length - 1} more steps (we already have the start node).

\textbf{Line 18:} Get neighbors of current node.

\textbf{Lines 19-20:} If no neighbors (dead end), stop the walk.

\textbf{Line 21:} \texttt{random.choice()} picks a random neighbor to move to.

\textbf{Line 22:} Add the new node to our walk.
\end{explanation}

\subsection{Node2Vec (2016)}

\begin{concept}
\textbf{Node2Vec} improves on DeepWalk with \textit{biased} random walks controlled by two parameters:
\begin{itemize}
    \item \textbf{p} (return parameter): Likelihood of returning to the previous node
    \item \textbf{q} (in-out parameter): Likelihood of exploring outward vs staying local
\end{itemize}

Low $p$: Encourage returning (local exploration)\\
Low $q$: Encourage moving outward (global exploration)
\end{concept}

\begin{mathbox}
\textbf{Node2Vec Transition Probabilities:}

From node $v$, having come from $t$, the unnormalized probability of going to node $x$ is:
\[
\alpha_{pq}(t, x) = \begin{cases}
\frac{1}{p} & \text{if } d_{tx} = 0 \text{ (return to } t\text{)} \\
1 & \text{if } d_{tx} = 1 \text{ (} x \text{ is neighbor of } t\text{)} \\
\frac{1}{q} & \text{if } d_{tx} = 2 \text{ (} x \text{ is not neighbor of } t\text{)}
\end{cases}
\]
where $d_{tx}$ is the shortest path distance from $t$ to $x$.
\end{mathbox}

\begin{lstlisting}[language=Python,caption={Node2Vec biased random walk}]
def node2vec_walk(adj_list, start_node, walk_length, p=1.0, q=1.0):
    """
    Generate a biased random walk (Node2Vec style).

    Args:
        adj_list: Dictionary mapping node to list of neighbors
        start_node: Node to start the walk from
        walk_length: Number of steps
        p: Return parameter (lower = more likely to return)
        q: In-out parameter (lower = more likely to go outward)

    Returns:
        List of nodes visited
    """
    walk = [start_node]

    if walk_length == 1:
        return walk

    # First step is uniform random
    neighbors = adj_list.get(start_node, [])
    if not neighbors:
        return walk
    walk.append(random.choice(neighbors))

    # Subsequent steps use biased sampling
    for _ in range(walk_length - 2):
        current = walk[-1]
        previous = walk[-2]
        neighbors = adj_list.get(current, [])

        if not neighbors:
            break

        # Calculate transition probabilities
        prev_neighbors = set(adj_list.get(previous, []))
        probabilities = []

        for neighbor in neighbors:
            if neighbor == previous:
                # Return to previous node
                probabilities.append(1.0 / p)
            elif neighbor in prev_neighbors:
                # Neighbor of previous (distance 1)
                probabilities.append(1.0)
            else:
                # Not neighbor of previous (distance 2)
                probabilities.append(1.0 / q)

        # Normalize probabilities
        prob_sum = sum(probabilities)
        probabilities = [p / prob_sum for p in probabilities]

        # Sample next node
        next_node = random.choices(neighbors, weights=probabilities, k=1)[0]
        walk.append(next_node)

    return walk

# Compare uniform vs biased walks
adj_list = {
    0: [1, 2, 3],
    1: [0, 4],
    2: [0, 5],
    3: [0, 6],
    4: [1, 7],
    5: [2, 7],
    6: [3, 7],
    7: [4, 5, 6]
}

print("Uniform walks (p=1, q=1):")
for _ in range(3):
    print(node2vec_walk(adj_list, 0, 6, p=1, q=1))

print("\nLocal walks (p=0.5, q=2):")
for _ in range(3):
    print(node2vec_walk(adj_list, 0, 6, p=0.5, q=2))

print("\nGlobal walks (p=2, q=0.5):")
for _ in range(3):
    print(node2vec_walk(adj_list, 0, 6, p=2, q=0.5))
\end{lstlisting}

\begin{explanation}
\textbf{Lines 27-28:} Keep track of current and previous nodes to compute transition probabilities.

\textbf{Line 35:} Get set of previous node's neighbors for fast lookup.

\textbf{Lines 38-47:} Assign probability based on Node2Vec rules---return, stay local, or explore.

\textbf{Lines 50-51:} Normalize so probabilities sum to 1.

\textbf{Line 54:} \texttt{random.choices()} samples with specified weights.
\end{explanation}

\begin{gitcommit}
\begin{lstlisting}[language=bash]
git add .
git commit -m "Implement random walk generators for DeepWalk and Node2Vec"
git push
\end{lstlisting}
\end{gitcommit}

%==============================================================================
% CHAPTER 6: GRAPH NEURAL NETWORKS
%==============================================================================
\section{Chapter 6: Graph Neural Networks (GNNs)}

\subsection{The Limitation of Random Walks}

Random walk methods don't use node features. In many applications, nodes have rich attributes (user profiles, document text, product descriptions). GNNs solve this.

\begin{concept}
\textbf{Graph Neural Networks} learn embeddings by iteratively aggregating information from neighbors. Each node's embedding is updated based on:
\begin{enumerate}
    \item Its own features
    \item Its neighbors' features
    \item The graph structure
\end{enumerate}
\end{concept}

\subsection{Message Passing Framework}

\begin{mathbox}
\textbf{Message Passing Neural Networks:}

For layer $l$, node $v$'s embedding is updated as:
\[
\mathbf{h}_v^{(l)} = \text{UPDATE}\left(\mathbf{h}_v^{(l-1)}, \text{AGGREGATE}\left(\left\{\mathbf{h}_u^{(l-1)} : u \in \mathcal{N}(v)\right\}\right)\right)
\]

Where:
\begin{itemize}
    \item $\mathbf{h}_v^{(l)}$ = embedding of node $v$ at layer $l$
    \item $\mathcal{N}(v)$ = neighbors of node $v$
    \item AGGREGATE = function to combine neighbor embeddings
    \item UPDATE = function to update node embedding
\end{itemize}
\end{mathbox}

\subsection{Graph Convolutional Network (GCN)}

\begin{mathbox}
\textbf{GCN Layer:}
\[
\mathbf{H}^{(l+1)} = \sigma\left(\tilde{D}^{-1/2}\tilde{A}\tilde{D}^{-1/2}\mathbf{H}^{(l)}\mathbf{W}^{(l)}\right)
\]

Where:
\begin{itemize}
    \item $\tilde{A} = A + I$ (adjacency with self-loops)
    \item $\tilde{D}$ = degree matrix of $\tilde{A}$
    \item $\mathbf{W}^{(l)}$ = learnable weight matrix
    \item $\sigma$ = activation function (e.g., ReLU)
\end{itemize}
\end{mathbox}

\begin{lstlisting}[language=Python,caption={GCN layer from scratch}]
import numpy as np

def gcn_layer(A, H, W):
    """
    Single GCN layer forward pass.

    Args:
        A: Adjacency matrix (n x n)
        H: Node features/embeddings (n x d_in)
        W: Weight matrix (d_in x d_out)

    Returns:
        Updated node embeddings (n x d_out)
    """
    n = A.shape[0]

    # Add self-loops: A_tilde = A + I
    A_tilde = A + np.eye(n)

    # Compute degree matrix D_tilde
    D_tilde = np.diag(np.sum(A_tilde, axis=1))

    # Compute D_tilde^(-1/2)
    D_tilde_inv_sqrt = np.diag(1.0 / np.sqrt(np.diag(D_tilde)))

    # Normalized adjacency: D^(-1/2) * A * D^(-1/2)
    A_norm = D_tilde_inv_sqrt @ A_tilde @ D_tilde_inv_sqrt

    # Message passing: aggregate neighbor features
    H_agg = A_norm @ H

    # Transform with learnable weights
    H_new = H_agg @ W

    # Apply activation (ReLU)
    H_new = np.maximum(0, H_new)

    return H_new

# Example: 4 nodes, each with 3 features, output 2 features
A = np.array([
    [0, 1, 1, 0],
    [1, 0, 1, 0],
    [1, 1, 0, 1],
    [0, 0, 1, 0],
], dtype=float)

H = np.array([
    [1.0, 0.0, 0.0],  # Node 0 features
    [0.0, 1.0, 0.0],  # Node 1 features
    [0.0, 0.0, 1.0],  # Node 2 features
    [1.0, 1.0, 0.0],  # Node 3 features
])

W = np.random.randn(3, 2) * 0.1  # 3 input features -> 2 output features

H_new = gcn_layer(A, H, W)
print(f"Input shape: {H.shape}")      # (4, 3)
print(f"Output shape: {H_new.shape}") # (4, 2)
print(f"\nNew embeddings:\n{H_new}")
\end{lstlisting}

\begin{explanation}
\textbf{Line 17:} Add self-loops so nodes include their own features during aggregation.

\textbf{Line 20:} Degree matrix: diagonal matrix where $D_{ii}$ = degree of node $i$.

\textbf{Line 23:} $D^{-1/2}$ for symmetric normalization---this prevents nodes with many neighbors from dominating.

\textbf{Line 26:} The normalized adjacency matrix. The normalization $D^{-1/2}AD^{-1/2}$ ensures messages are scaled by node degrees.

\textbf{Line 29:} Matrix multiplication $A_{norm} \cdot H$ aggregates neighbor features. Each node gets a weighted sum of its neighbors' features.

\textbf{Line 32:} Linear transformation with learnable weights.

\textbf{Line 35:} ReLU activation: $\max(0, x)$ introduces non-linearity.
\end{explanation}

\begin{warning}
The manual GCN implementation above is for understanding. In practice, always use libraries like PyTorch Geometric which handle:
\begin{itemize}
    \item GPU acceleration
    \item Sparse matrix operations
    \item Batching for large graphs
    \item Numerical stability
\end{itemize}
\end{warning}

\subsection{GraphSAGE: Scalable GNNs}

\begin{concept}
\textbf{GraphSAGE} (SAmple and aggreGatE) solves GCN's scalability issues:
\begin{enumerate}
    \item Instead of using all neighbors, \textit{sample} a fixed number
    \item Use different aggregation functions (mean, max, LSTM)
    \item Can generate embeddings for unseen nodes (inductive learning)
\end{enumerate}
\end{concept}

\begin{mathbox}
\textbf{GraphSAGE Update:}
\[
\mathbf{h}_v^{(l)} = \sigma\left(\mathbf{W}^{(l)} \cdot \text{CONCAT}\left(\mathbf{h}_v^{(l-1)}, \text{AGG}\left(\left\{\mathbf{h}_u^{(l-1)} : u \in \mathcal{N}_s(v)\right\}\right)\right)\right)
\]
Where $\mathcal{N}_s(v)$ is a sampled subset of neighbors.
\end{mathbox}

\begin{lstlisting}[language=Python,caption={GraphSAGE mean aggregation}]
def graphsage_layer(adj_list, H, W, sample_size=3):
    """
    GraphSAGE layer with mean aggregation.

    Args:
        adj_list: Adjacency list
        H: Node embeddings (n x d_in)
        W: Weight matrix (2*d_in x d_out)
        sample_size: Number of neighbors to sample

    Returns:
        Updated embeddings (n x d_out)
    """
    n = H.shape[0]
    d_in = H.shape[1]
    d_out = W.shape[1]

    H_new = np.zeros((n, d_out))

    for v in range(n):
        # Sample neighbors
        neighbors = adj_list.get(v, [])
        if len(neighbors) > sample_size:
            sampled = random.sample(neighbors, sample_size)
        else:
            sampled = neighbors

        # Aggregate neighbor embeddings (mean)
        if sampled:
            neighbor_embs = H[sampled]
            agg = np.mean(neighbor_embs, axis=0)
        else:
            agg = np.zeros(d_in)

        # Concatenate self and aggregated neighbors
        combined = np.concatenate([H[v], agg])

        # Transform and activate
        H_new[v] = np.maximum(0, combined @ W)

    return H_new
\end{lstlisting}

\begin{explanation}
\textbf{Lines 22-26:} Sample a fixed number of neighbors. This makes computation time constant per node regardless of degree.

\textbf{Lines 29-33:} Mean aggregation: average the neighbor embeddings.

\textbf{Line 36:} Concatenate node's own embedding with aggregated neighbor embedding.

\textbf{Line 39:} Transform with weights and apply ReLU.
\end{explanation}

\subsection{Graph Attention Networks (GAT)}

\begin{concept}
\textbf{GAT} learns \textit{attention weights} to decide how much to weight each neighbor. Some neighbors are more important than others---GAT learns this automatically.
\end{concept}

\begin{mathbox}
\textbf{GAT Attention:}
\[
\alpha_{ij} = \frac{\exp\left(\text{LeakyReLU}\left(\mathbf{a}^T [\mathbf{W}\mathbf{h}_i || \mathbf{W}\mathbf{h}_j]\right)\right)}{\sum_{k \in \mathcal{N}(i)} \exp\left(\text{LeakyReLU}\left(\mathbf{a}^T [\mathbf{W}\mathbf{h}_i || \mathbf{W}\mathbf{h}_k]\right)\right)}
\]
\[
\mathbf{h}_i' = \sigma\left(\sum_{j \in \mathcal{N}(i)} \alpha_{ij} \mathbf{W}\mathbf{h}_j\right)
\]
\end{mathbox}

\begin{lstlisting}[language=Python,caption={Simplified GAT attention}]
def compute_attention(h_i, h_j, W, a):
    """
    Compute attention coefficient between nodes i and j.

    Args:
        h_i: Source node embedding
        h_j: Target node embedding
        W: Weight matrix
        a: Attention vector

    Returns:
        Unnormalized attention score
    """
    # Transform embeddings
    Wh_i = W @ h_i
    Wh_j = W @ h_j

    # Concatenate
    concat = np.concatenate([Wh_i, Wh_j])

    # Compute attention score
    score = a @ concat

    # LeakyReLU
    score = score if score > 0 else 0.01 * score

    return score

def softmax(x):
    """Numerically stable softmax."""
    e_x = np.exp(x - np.max(x))
    return e_x / e_x.sum()

def gat_aggregate(v, adj_list, H, W, a):
    """
    GAT aggregation for node v.

    Args:
        v: Node index
        adj_list: Adjacency list
        H: All node embeddings
        W: Weight matrix
        a: Attention vector

    Returns:
        New embedding for node v
    """
    neighbors = adj_list.get(v, [])
    if not neighbors:
        return W @ H[v]

    # Compute attention scores
    scores = []
    for u in neighbors:
        score = compute_attention(H[v], H[u], W, a)
        scores.append(score)

    # Normalize with softmax
    attention_weights = softmax(np.array(scores))

    # Weighted aggregation
    agg = np.zeros_like(W @ H[v])
    for u, alpha in zip(neighbors, attention_weights):
        agg += alpha * (W @ H[u])

    return agg
\end{lstlisting}

\begin{checkpoint}
Verify you understand message passing:
\begin{lstlisting}[language=Python]
# For a node with neighbors [1, 2, 3] with features
# h1 = [1, 0], h2 = [0, 1], h3 = [1, 1]
# What is the mean aggregation?
neighbors_h = np.array([[1, 0], [0, 1], [1, 1]])
mean_agg = np.mean(neighbors_h, axis=0)
print(mean_agg)  # Should be [0.667, 0.667]
\end{lstlisting}
\end{checkpoint}

\begin{gitcommit}
\begin{lstlisting}[language=bash]
git add .
git commit -m "Implement GCN, GraphSAGE, and GAT layers from scratch"
git push
\end{lstlisting}
\end{gitcommit}

%==============================================================================
% PART III: BUILDING THE RETRIEVAL SYSTEM
%==============================================================================
\part{Building the Retrieval System}

\section{Chapter 7: Setting Up the Environment}

\subsection{Step 7.1: Create Project Structure}

\begin{lstlisting}[language=bash,caption={Create project}]
mkdir graph-retrieval
cd graph-retrieval
python -m venv venv
source venv/bin/activate  # On Windows: venv\Scripts\activate
\end{lstlisting}

\begin{explanation}
\textbf{Line 3:} Create a virtual environment---an isolated Python installation for this project.

\textbf{Line 4:} Activate the environment. All subsequent \texttt{pip install} commands will install to this environment.
\end{explanation}

\subsection{Step 7.2: Install Dependencies}

\begin{lstlisting}[language=bash,caption={Install packages}]
pip install torch torchvision
pip install torch-geometric
pip install numpy pandas scikit-learn
pip install faiss-cpu  # For fast similarity search
pip install networkx   # For graph operations
pip install matplotlib # For visualization
\end{lstlisting}

\begin{explanation}
\textbf{Line 1:} PyTorch - the deep learning framework.

\textbf{Line 2:} PyTorch Geometric - GNN library built on PyTorch.

\textbf{Line 4:} FAISS - Facebook's library for fast similarity search in high-dimensional spaces.

\textbf{Line 5:} NetworkX - Python library for graph analysis.
\end{explanation}

\begin{checkpoint}
Verify installation:
\begin{lstlisting}[language=Python]
import torch
import torch_geometric
import faiss
print(f"PyTorch: {torch.__version__}")
print(f"PyG: {torch_geometric.__version__}")
print("FAISS: OK")
\end{lstlisting}
\end{checkpoint}

\subsection{Step 7.3: Project Structure}

\begin{lstlisting}[language=bash,caption={Directory structure}]
graph-retrieval/
|-- data/
|   |-- raw/            # Raw datasets
|   |-- processed/      # Processed data
|-- models/
|   |-- gnn.py          # GNN model definitions
|   |-- encoder.py      # Embedding encoder
|-- retrieval/
|   |-- index.py        # FAISS index management
|   |-- search.py       # Search functionality
|-- utils/
|   |-- data_loader.py  # Data loading utilities
|   |-- metrics.py      # Evaluation metrics
|-- train.py            # Training script
|-- evaluate.py         # Evaluation script
|-- config.py           # Configuration
|-- requirements.txt    # Dependencies
\end{lstlisting}

\begin{gitcommit}
\begin{lstlisting}[language=bash]
git init
git add .
git commit -m "Initial project structure and dependencies"
\end{lstlisting}
\end{gitcommit}

%==============================================================================
% CHAPTER 8: IMPLEMENTING THE GNN
%==============================================================================
\section{Chapter 8: Implementing the GNN Model}

\subsection{Step 8.1: Create the GNN Encoder}

Create \texttt{models/gnn.py}:

\begin{lstlisting}[language=Python,caption={models/gnn.py}]
import torch
import torch.nn as nn
import torch.nn.functional as F
from torch_geometric.nn import GCNConv, SAGEConv, GATConv

class GraphEncoder(nn.Module):
    """
    Graph Neural Network encoder for generating node embeddings.

    This module takes a graph (node features + edge structure) and
    produces dense embeddings for each node that capture both
    the node's features and its structural position in the graph.
    """

    def __init__(
        self,
        input_dim: int,
        hidden_dim: int,
        output_dim: int,
        num_layers: int = 2,
        dropout: float = 0.5,
        gnn_type: str = 'gcn'
    ):
        """
        Initialize the graph encoder.

        Args:
            input_dim: Dimension of input node features
            hidden_dim: Dimension of hidden layers
            output_dim: Dimension of output embeddings
            num_layers: Number of GNN layers
            dropout: Dropout probability
            gnn_type: Type of GNN ('gcn', 'sage', or 'gat')
        """
        super(GraphEncoder, self).__init__()

        self.num_layers = num_layers
        self.dropout = dropout

        # Create GNN layers
        self.convs = nn.ModuleList()

        # Select GNN layer type
        if gnn_type == 'gcn':
            ConvLayer = GCNConv
        elif gnn_type == 'sage':
            ConvLayer = SAGEConv
        elif gnn_type == 'gat':
            ConvLayer = lambda in_c, out_c: GATConv(in_c, out_c, heads=4, concat=False)
        else:
            raise ValueError(f"Unknown GNN type: {gnn_type}")

        # First layer: input_dim -> hidden_dim
        self.convs.append(ConvLayer(input_dim, hidden_dim))

        # Middle layers: hidden_dim -> hidden_dim
        for _ in range(num_layers - 2):
            self.convs.append(ConvLayer(hidden_dim, hidden_dim))

        # Last layer: hidden_dim -> output_dim
        if num_layers > 1:
            self.convs.append(ConvLayer(hidden_dim, output_dim))

    def forward(self, x, edge_index):
        """
        Forward pass through the GNN.

        Args:
            x: Node feature matrix [num_nodes, input_dim]
            edge_index: Edge indices [2, num_edges]

        Returns:
            Node embeddings [num_nodes, output_dim]
        """
        for i, conv in enumerate(self.convs):
            x = conv(x, edge_index)

            # Apply activation and dropout to all but last layer
            if i < self.num_layers - 1:
                x = F.relu(x)
                x = F.dropout(x, p=self.dropout, training=self.training)

        # L2 normalize embeddings (important for cosine similarity)
        x = F.normalize(x, p=2, dim=1)

        return x


class DualEncoder(nn.Module):
    """
    Dual encoder for query and document embedding.

    Uses separate encoders for queries and documents, allowing
    different representations while mapping to the same space.
    """

    def __init__(
        self,
        input_dim: int,
        hidden_dim: int,
        output_dim: int,
        num_layers: int = 2,
        share_weights: bool = True
    ):
        super(DualEncoder, self).__init__()

        self.query_encoder = GraphEncoder(
            input_dim, hidden_dim, output_dim, num_layers
        )

        if share_weights:
            self.doc_encoder = self.query_encoder
        else:
            self.doc_encoder = GraphEncoder(
                input_dim, hidden_dim, output_dim, num_layers
            )

    def encode_query(self, x, edge_index):
        return self.query_encoder(x, edge_index)

    def encode_document(self, x, edge_index):
        return self.doc_encoder(x, edge_index)

    def forward(self, query_x, query_edge, doc_x, doc_edge):
        query_emb = self.encode_query(query_x, query_edge)
        doc_emb = self.encode_document(doc_x, doc_edge)
        return query_emb, doc_emb
\end{lstlisting}

\begin{explanation}
\textbf{Lines 6-13:} Class docstring explaining the purpose---GNNs produce embeddings capturing features AND structure.

\textbf{Line 34:} \texttt{super().\_\_init\_\_()} initializes the parent \texttt{nn.Module} class.

\textbf{Line 40:} \texttt{nn.ModuleList} is a container that properly registers layers for PyTorch.

\textbf{Lines 42-48:} Select the GNN layer type. The lambda for GAT configures 4 attention heads.

\textbf{Lines 51-58:} Stack multiple GNN layers. First layer handles input dimension, middle layers maintain hidden dimension, last layer produces output dimension.

\textbf{Lines 73-77:} Apply ReLU activation and dropout between layers, but NOT after the last layer.

\textbf{Line 80:} L2 normalization ensures embeddings have unit length---this is crucial for cosine similarity to work correctly.

\textbf{Lines 84-112:} Dual encoder pattern---common in retrieval where queries and documents might need different processing.
\end{explanation}

\begin{concept}
\textbf{Why L2 Normalize Embeddings?}

After L2 normalization, $||\mathbf{x}|| = 1$ for all embeddings. This means:
\[
\text{cosine\_sim}(\mathbf{a}, \mathbf{b}) = \mathbf{a} \cdot \mathbf{b}
\]
The dot product becomes the cosine similarity, making similarity computation extremely fast.
\end{concept}

\subsection{Step 8.2: Implement the Training Loop}

Create \texttt{train.py}:

\begin{lstlisting}[language=Python,caption={train.py}]
import torch
import torch.nn as nn
import torch.optim as optim
from torch_geometric.data import Data, DataLoader
import numpy as np
from models.gnn import GraphEncoder

def contrastive_loss(query_emb, pos_emb, neg_emb, margin=0.2):
    """
    Triplet margin loss for contrastive learning.

    Encourages: sim(query, positive) > sim(query, negative) + margin

    Args:
        query_emb: Query embeddings [batch_size, dim]
        pos_emb: Positive (relevant) embeddings [batch_size, dim]
        neg_emb: Negative (irrelevant) embeddings [batch_size, dim]
        margin: Minimum margin between positive and negative

    Returns:
        Loss value
    """
    # Cosine similarity (embeddings are L2 normalized)
    pos_sim = torch.sum(query_emb * pos_emb, dim=1)
    neg_sim = torch.sum(query_emb * neg_emb, dim=1)

    # Triplet loss: max(0, margin - pos_sim + neg_sim)
    loss = torch.clamp(margin - pos_sim + neg_sim, min=0)

    return loss.mean()


def info_nce_loss(query_emb, doc_emb, temperature=0.07):
    """
    InfoNCE loss (used in CLIP, SimCLR, etc.)

    Treats retrieval as a classification problem:
    Given query i, classify which document is the match.

    Args:
        query_emb: Query embeddings [batch_size, dim]
        doc_emb: Document embeddings [batch_size, dim]
        temperature: Softmax temperature (lower = sharper)

    Returns:
        Loss value
    """
    # Similarity matrix [batch_size, batch_size]
    logits = torch.matmul(query_emb, doc_emb.T) / temperature

    # Labels: diagonal entries are positives (query i matches doc i)
    labels = torch.arange(logits.size(0), device=logits.device)

    # Cross entropy loss
    loss = nn.CrossEntropyLoss()(logits, labels)

    return loss


def train_epoch(model, data_loader, optimizer, device):
    """
    Train for one epoch.

    Args:
        model: The GNN model
        data_loader: DataLoader yielding batches
        optimizer: Optimizer
        device: Device to train on

    Returns:
        Average loss for the epoch
    """
    model.train()
    total_loss = 0

    for batch in data_loader:
        batch = batch.to(device)
        optimizer.zero_grad()

        # Forward pass
        embeddings = model(batch.x, batch.edge_index)

        # Compute loss (using InfoNCE)
        # Assuming batch contains paired query-document data
        batch_size = embeddings.size(0) // 2
        query_emb = embeddings[:batch_size]
        doc_emb = embeddings[batch_size:]

        loss = info_nce_loss(query_emb, doc_emb)

        # Backward pass
        loss.backward()
        optimizer.step()

        total_loss += loss.item()

    return total_loss / len(data_loader)


def train(
    model,
    train_loader,
    val_loader,
    num_epochs=100,
    learning_rate=0.001,
    device='cuda'
):
    """
    Full training loop with validation.

    Args:
        model: The GNN model
        train_loader: Training data loader
        val_loader: Validation data loader
        num_epochs: Number of training epochs
        learning_rate: Learning rate
        device: Device to train on

    Returns:
        Trained model
    """
    model = model.to(device)
    optimizer = optim.Adam(model.parameters(), lr=learning_rate)
    scheduler = optim.lr_scheduler.ReduceLROnPlateau(
        optimizer, mode='min', patience=10, factor=0.5
    )

    best_val_loss = float('inf')

    for epoch in range(num_epochs):
        # Training
        train_loss = train_epoch(model, train_loader, optimizer, device)

        # Validation
        model.eval()
        val_loss = 0
        with torch.no_grad():
            for batch in val_loader:
                batch = batch.to(device)
                embeddings = model(batch.x, batch.edge_index)
                batch_size = embeddings.size(0) // 2
                query_emb = embeddings[:batch_size]
                doc_emb = embeddings[batch_size:]
                val_loss += info_nce_loss(query_emb, doc_emb).item()

        val_loss /= len(val_loader)
        scheduler.step(val_loss)

        # Logging
        print(f"Epoch {epoch+1}/{num_epochs}")
        print(f"  Train Loss: {train_loss:.4f}")
        print(f"  Val Loss: {val_loss:.4f}")
        print(f"  LR: {optimizer.param_groups[0]['lr']:.6f}")

        # Save best model
        if val_loss < best_val_loss:
            best_val_loss = val_loss
            torch.save(model.state_dict(), 'best_model.pt')
            print("  Saved new best model!")

    # Load best model
    model.load_state_dict(torch.load('best_model.pt'))
    return model
\end{lstlisting}

\begin{explanation}
\textbf{Lines 8-29:} \textbf{Triplet loss}: Classic metric learning loss. Push positive examples closer, negative examples farther.

\textbf{Lines 32-55:} \textbf{InfoNCE loss}: More modern approach used in CLIP. Treats retrieval as classification---given a query, which document in the batch is correct?

\textbf{Line 48:} Divide by temperature. Lower temperature makes the softmax sharper (more confident predictions).

\textbf{Line 51:} The labels are 0, 1, 2, ..., batch\_size-1 because query $i$ should match document $i$ (diagonal of similarity matrix).

\textbf{Lines 111-113:} Learning rate scheduler: if validation loss doesn't improve for 10 epochs, reduce LR by half.

\textbf{Lines 140-142:} Save the best model based on validation loss. This prevents overfitting.
\end{explanation}

\begin{warning}
\textbf{Common Training Mistake:}

Always use \texttt{model.eval()} and \texttt{torch.no\_grad()} during validation/inference. Otherwise:
\begin{itemize}
    \item Dropout will still be active (wrong predictions)
    \item Gradients will be computed (wasted memory)
\end{itemize}
\end{warning}

\begin{checkpoint}
Create a dummy test:
\begin{lstlisting}[language=Python]
import torch
from models.gnn import GraphEncoder

model = GraphEncoder(input_dim=16, hidden_dim=32, output_dim=64)
x = torch.randn(100, 16)  # 100 nodes, 16 features
edge_index = torch.randint(0, 100, (2, 500))  # 500 random edges

out = model(x, edge_index)
print(f"Output shape: {out.shape}")  # Should be [100, 64]
print(f"Embeddings normalized: {torch.allclose(out.norm(dim=1), torch.ones(100))}")
\end{lstlisting}
\end{checkpoint}

\begin{gitcommit}
\begin{lstlisting}[language=bash]
git add .
git commit -m "Implement GNN encoder and training loop with contrastive losses"
git push
\end{lstlisting}
\end{gitcommit}

%==============================================================================
% CHAPTER 9: BUILDING THE RETRIEVAL INDEX
%==============================================================================
\section{Chapter 9: Building the Retrieval Index}

\subsection{Step 9.1: FAISS Index for Fast Search}

\begin{concept}
\textbf{The Search Problem:}

Given a query embedding and millions of document embeddings, how do we find the most similar ones quickly?

Naive approach: Compute similarity with every document. $O(n)$ per query---too slow!

\textbf{Solution}: Use approximate nearest neighbor (ANN) search with FAISS.
\end{concept}

Create \texttt{retrieval/index.py}:

\begin{lstlisting}[language=Python,caption={retrieval/index.py}]
import numpy as np
import faiss

class FaissIndex:
    """
    FAISS-based index for fast approximate nearest neighbor search.

    Supports both exact search (brute force) and approximate search
    (using IVF quantization for large-scale retrieval).
    """

    def __init__(self, dimension: int, index_type: str = 'flat'):
        """
        Initialize the FAISS index.

        Args:
            dimension: Embedding dimension
            index_type: Type of index
                - 'flat': Exact search (brute force)
                - 'ivf': Approximate search (faster for large datasets)
                - 'hnsw': Graph-based approximate search
        """
        self.dimension = dimension
        self.index_type = index_type
        self.index = None
        self.id_map = {}  # Maps FAISS internal ID to external ID

        self._build_index()

    def _build_index(self):
        """Build the FAISS index based on type."""
        if self.index_type == 'flat':
            # Exact search with cosine similarity
            # (FAISS uses inner product, so we normalize vectors)
            self.index = faiss.IndexFlatIP(self.dimension)

        elif self.index_type == 'ivf':
            # IVF index for approximate search
            # Partitions space into clusters, searches only nearby clusters
            nlist = 100  # Number of clusters
            quantizer = faiss.IndexFlatIP(self.dimension)
            self.index = faiss.IndexIVFFlat(
                quantizer, self.dimension, nlist, faiss.METRIC_INNER_PRODUCT
            )

        elif self.index_type == 'hnsw':
            # HNSW: Hierarchical Navigable Small World graphs
            # Very fast for high-recall approximate search
            M = 32  # Number of connections per layer
            self.index = faiss.IndexHNSWFlat(self.dimension, M)
            self.index.metric_type = faiss.METRIC_INNER_PRODUCT

        else:
            raise ValueError(f"Unknown index type: {self.index_type}")

    def train(self, embeddings: np.ndarray):
        """
        Train the index (required for IVF).

        Args:
            embeddings: Training embeddings [n_samples, dimension]
        """
        if self.index_type == 'ivf':
            if not self.index.is_trained:
                self.index.train(embeddings.astype(np.float32))

    def add(self, embeddings: np.ndarray, ids: list = None):
        """
        Add embeddings to the index.

        Args:
            embeddings: Embeddings to add [n_samples, dimension]
            ids: Optional external IDs for each embedding
        """
        n = embeddings.shape[0]

        # Normalize for cosine similarity
        embeddings = embeddings.astype(np.float32)
        faiss.normalize_L2(embeddings)

        # Map IDs
        start_id = len(self.id_map)
        if ids is None:
            ids = list(range(start_id, start_id + n))

        for i, ext_id in enumerate(ids):
            self.id_map[start_id + i] = ext_id

        # Add to index
        self.index.add(embeddings)

    def search(self, query: np.ndarray, k: int = 10):
        """
        Search for k nearest neighbors.

        Args:
            query: Query embedding(s) [n_queries, dimension]
            k: Number of results to return

        Returns:
            Tuple of (distances, ids) where:
                - distances: [n_queries, k] similarity scores
                - ids: [n_queries, k] external IDs of results
        """
        query = query.astype(np.float32)
        if query.ndim == 1:
            query = query.reshape(1, -1)

        # Normalize query
        faiss.normalize_L2(query)

        # Search
        distances, indices = self.index.search(query, k)

        # Map internal IDs to external IDs
        external_ids = np.array([
            [self.id_map.get(int(idx), -1) for idx in row]
            for row in indices
        ])

        return distances, external_ids

    def save(self, path: str):
        """Save index to disk."""
        faiss.write_index(self.index, path + '.index')
        np.save(path + '.id_map.npy', self.id_map)

    def load(self, path: str):
        """Load index from disk."""
        self.index = faiss.read_index(path + '.index')
        self.id_map = np.load(path + '.id_map.npy', allow_pickle=True).item()
\end{lstlisting}

\begin{explanation}
\textbf{Line 34:} \texttt{IndexFlatIP} = Flat index with Inner Product. Since we normalize vectors, inner product equals cosine similarity.

\textbf{Lines 38-42:} \textbf{IVF (Inverted File Index)}: Clusters the data, then searches only nearby clusters. Much faster for large datasets but approximate.

\textbf{Lines 45-48:} \textbf{HNSW}: Creates a graph structure for navigation. State-of-the-art for recall/speed tradeoff.

\textbf{Lines 73-74:} \texttt{faiss.normalize\_L2()} normalizes in-place for efficiency.

\textbf{Lines 76-81:} ID mapping: FAISS uses sequential internal IDs, but we want custom external IDs (like document IDs).
\end{explanation}

\begin{mathbox}
\textbf{ANN Search Complexity:}

\begin{tabular}{|l|l|l|}
\hline
\textbf{Method} & \textbf{Build} & \textbf{Query} \\
\hline
Flat (Brute Force) & $O(n)$ & $O(n)$ \\
IVF & $O(n\sqrt{n})$ & $O(\sqrt{n})$ \\
HNSW & $O(n\log n)$ & $O(\log n)$ \\
\hline
\end{tabular}

For 1 million documents, HNSW can be 1000x faster than brute force!
\end{mathbox}

\subsection{Step 9.2: The Retrieval Pipeline}

Create \texttt{retrieval/search.py}:

\begin{lstlisting}[language=Python,caption={retrieval/search.py}]
import torch
import numpy as np
from typing import List, Dict, Any

class GraphRetriever:
    """
    Complete retrieval pipeline using graph embeddings.

    This class combines:
    1. GNN encoder for generating embeddings
    2. FAISS index for fast similarity search
    3. Post-processing for result formatting
    """

    def __init__(self, model, index, device='cuda'):
        """
        Initialize the retriever.

        Args:
            model: Trained GNN encoder
            index: FAISS index with document embeddings
            device: Device for model inference
        """
        self.model = model.to(device)
        self.model.eval()
        self.index = index
        self.device = device

    @torch.no_grad()
    def encode(self, x, edge_index):
        """
        Encode a query graph to embedding.

        Args:
            x: Node features [num_nodes, feature_dim]
            edge_index: Edge indices [2, num_edges]

        Returns:
            Embedding vector [embedding_dim]
        """
        x = torch.FloatTensor(x).to(self.device)
        edge_index = torch.LongTensor(edge_index).to(self.device)

        embedding = self.model(x, edge_index)

        # If multiple nodes, aggregate (mean pooling)
        if embedding.shape[0] > 1:
            embedding = embedding.mean(dim=0, keepdim=True)

        return embedding.cpu().numpy().flatten()

    def search(
        self,
        query_features,
        query_edges,
        k: int = 10,
        filter_fn=None
    ) -> List[Dict[str, Any]]:
        """
        Search for relevant documents given a query.

        Args:
            query_features: Query node features
            query_edges: Query edge indices
            k: Number of results to return
            filter_fn: Optional function to filter results

        Returns:
            List of results with ids, scores, and metadata
        """
        # Encode query
        query_embedding = self.encode(query_features, query_edges)

        # Search index
        scores, ids = self.index.search(
            query_embedding.reshape(1, -1), k * 2  # Over-fetch for filtering
        )

        # Format results
        results = []
        for score, doc_id in zip(scores[0], ids[0]):
            if doc_id == -1:
                continue

            result = {
                'id': int(doc_id),
                'score': float(score),
            }

            # Apply filter if provided
            if filter_fn and not filter_fn(result):
                continue

            results.append(result)

            if len(results) >= k:
                break

        return results

    def batch_search(
        self,
        queries: List[tuple],
        k: int = 10
    ) -> List[List[Dict[str, Any]]]:
        """
        Search for multiple queries at once.

        Args:
            queries: List of (features, edges) tuples
            k: Number of results per query

        Returns:
            List of result lists
        """
        # Encode all queries
        query_embeddings = np.stack([
            self.encode(features, edges)
            for features, edges in queries
        ])

        # Batch search
        scores, ids = self.index.search(query_embeddings, k)

        # Format results
        all_results = []
        for query_scores, query_ids in zip(scores, ids):
            results = [
                {'id': int(doc_id), 'score': float(score)}
                for score, doc_id in zip(query_scores, query_ids)
                if doc_id != -1
            ]
            all_results.append(results)

        return all_results
\end{lstlisting}

\begin{explanation}
\textbf{Line 28:} \texttt{@torch.no\_grad()} decorator disables gradient computation for efficiency.

\textbf{Lines 44-46:} Graph pooling: If the query graph has multiple nodes, we average their embeddings to get a single query vector.

\textbf{Line 70:} Over-fetch results to allow for filtering.

\textbf{Lines 97-105:} Batch encoding is more efficient than encoding one at a time due to GPU parallelism.
\end{explanation}

\begin{checkpoint}
Test the retrieval pipeline:
\begin{lstlisting}[language=Python]
import numpy as np
from retrieval.index import FaissIndex

# Create and populate index
index = FaissIndex(dimension=64)
docs = np.random.randn(1000, 64)  # 1000 documents
index.add(docs)

# Search
query = np.random.randn(1, 64)
scores, ids = index.search(query, k=5)
print(f"Top 5 results: {ids[0]}")
print(f"Scores: {scores[0]}")
\end{lstlisting}
\end{checkpoint}

\begin{gitcommit}
\begin{lstlisting}[language=bash]
git add .
git commit -m "Implement FAISS index and retrieval pipeline"
git push
\end{lstlisting}
\end{gitcommit}

%==============================================================================
% CHAPTER 10: EVALUATION
%==============================================================================
\section{Chapter 10: Evaluation Metrics}

\subsection{Retrieval Metrics}

Create \texttt{utils/metrics.py}:

\begin{lstlisting}[language=Python,caption={utils/metrics.py}]
import numpy as np
from typing import List, Set

def precision_at_k(retrieved: List[int], relevant: Set[int], k: int) -> float:
    """
    Precision@K: What fraction of top-K results are relevant?

    Args:
        retrieved: List of retrieved document IDs (ranked)
        relevant: Set of relevant document IDs
        k: Number of results to consider

    Returns:
        Precision@K score (0 to 1)
    """
    if k == 0:
        return 0.0

    retrieved_k = retrieved[:k]
    relevant_retrieved = sum(1 for doc_id in retrieved_k if doc_id in relevant)

    return relevant_retrieved / k


def recall_at_k(retrieved: List[int], relevant: Set[int], k: int) -> float:
    """
    Recall@K: What fraction of relevant documents are in top-K?

    Args:
        retrieved: List of retrieved document IDs (ranked)
        relevant: Set of relevant document IDs
        k: Number of results to consider

    Returns:
        Recall@K score (0 to 1)
    """
    if len(relevant) == 0:
        return 0.0

    retrieved_k = retrieved[:k]
    relevant_retrieved = sum(1 for doc_id in retrieved_k if doc_id in relevant)

    return relevant_retrieved / len(relevant)


def average_precision(retrieved: List[int], relevant: Set[int]) -> float:
    """
    Average Precision (AP): Area under precision-recall curve.

    Higher AP means relevant documents are ranked higher.

    Args:
        retrieved: List of retrieved document IDs (ranked)
        relevant: Set of relevant document IDs

    Returns:
        AP score (0 to 1)
    """
    if len(relevant) == 0:
        return 0.0

    precisions = []
    num_relevant_seen = 0

    for i, doc_id in enumerate(retrieved):
        if doc_id in relevant:
            num_relevant_seen += 1
            precision = num_relevant_seen / (i + 1)
            precisions.append(precision)

    if not precisions:
        return 0.0

    return sum(precisions) / len(relevant)


def mean_average_precision(
    all_retrieved: List[List[int]],
    all_relevant: List[Set[int]]
) -> float:
    """
    Mean Average Precision (MAP): Average AP across all queries.

    The standard metric for ranking quality.

    Args:
        all_retrieved: List of retrieved lists, one per query
        all_relevant: List of relevant sets, one per query

    Returns:
        MAP score (0 to 1)
    """
    aps = [
        average_precision(retrieved, relevant)
        for retrieved, relevant in zip(all_retrieved, all_relevant)
    ]

    return np.mean(aps) if aps else 0.0


def ndcg_at_k(
    retrieved: List[int],
    relevance_scores: dict,
    k: int
) -> float:
    """
    Normalized Discounted Cumulative Gain (NDCG@K).

    Unlike precision/recall, NDCG handles graded relevance
    (e.g., highly relevant, somewhat relevant, not relevant).

    Args:
        retrieved: List of retrieved document IDs (ranked)
        relevance_scores: Dict mapping doc_id to relevance score
        k: Number of results to consider

    Returns:
        NDCG@K score (0 to 1)
    """
    def dcg(scores):
        return sum(
            score / np.log2(i + 2)  # +2 because log2(1) = 0
            for i, score in enumerate(scores)
        )

    # Get relevance scores for retrieved docs
    retrieved_k = retrieved[:k]
    retrieved_scores = [relevance_scores.get(doc_id, 0) for doc_id in retrieved_k]

    # Calculate DCG
    actual_dcg = dcg(retrieved_scores)

    # Calculate ideal DCG (perfect ranking)
    ideal_scores = sorted(relevance_scores.values(), reverse=True)[:k]
    ideal_dcg = dcg(ideal_scores)

    if ideal_dcg == 0:
        return 0.0

    return actual_dcg / ideal_dcg


def mrr(all_retrieved: List[List[int]], all_relevant: List[Set[int]]) -> float:
    """
    Mean Reciprocal Rank (MRR).

    Focuses on where the first relevant result appears.
    MRR = average of 1/rank_of_first_relevant.

    Args:
        all_retrieved: List of retrieved lists, one per query
        all_relevant: List of relevant sets, one per query

    Returns:
        MRR score (0 to 1)
    """
    reciprocal_ranks = []

    for retrieved, relevant in zip(all_retrieved, all_relevant):
        for i, doc_id in enumerate(retrieved):
            if doc_id in relevant:
                reciprocal_ranks.append(1.0 / (i + 1))
                break
        else:
            reciprocal_ranks.append(0.0)

    return np.mean(reciprocal_ranks) if reciprocal_ranks else 0.0
\end{lstlisting}

\begin{mathbox}
\textbf{Key Metrics Summary:}

\begin{itemize}
    \item \textbf{Precision@K}: $\frac{\text{relevant in top K}}{K}$
    \item \textbf{Recall@K}: $\frac{\text{relevant in top K}}{\text{total relevant}}$
    \item \textbf{MAP}: Average precision across all relevant documents
    \item \textbf{NDCG}: Handles graded relevance with position discounting
    \item \textbf{MRR}: Focuses on rank of first relevant result
\end{itemize}
\end{mathbox}

\begin{concept}
\textbf{When to Use Each Metric:}

\begin{itemize}
    \item \textbf{Precision@K}: When you show exactly K results (e.g., top 10)
    \item \textbf{Recall@K}: When finding all relevant docs matters
    \item \textbf{MAP}: General ranking quality
    \item \textbf{NDCG}: When relevance is graded (stars, ratings)
    \item \textbf{MRR}: When users typically click first relevant result
\end{itemize}
\end{concept}

\begin{checkpoint}
Test metrics:
\begin{lstlisting}[language=Python]
from utils.metrics import precision_at_k, recall_at_k, average_precision

retrieved = [1, 5, 3, 8, 2]  # Ranked results
relevant = {1, 2, 3}  # Ground truth

print(f"P@3: {precision_at_k(retrieved, relevant, 3):.3f}")  # 2/3 = 0.667
print(f"R@3: {recall_at_k(retrieved, relevant, 3):.3f}")    # 2/3 = 0.667
print(f"AP: {average_precision(retrieved, relevant):.3f}")  # (1/1 + 2/3 + 3/5) / 3
\end{lstlisting}
\end{checkpoint}

\begin{gitcommit}
\begin{lstlisting}[language=bash]
git add .
git commit -m "Implement retrieval evaluation metrics"
git push
\end{lstlisting}
\end{gitcommit}

%==============================================================================
% CHAPTER 11: PUTTING IT ALL TOGETHER
%==============================================================================
\section{Chapter 11: Complete Example}

\subsection{Step 11.1: End-to-End Pipeline}

Create \texttt{example.py}:

\begin{lstlisting}[language=Python,caption={example.py - Complete pipeline}]
"""
Complete example: Build a retrieval system for a citation network.

Documents are papers, edges are citations.
Query: Find papers similar to a given paper.
"""

import torch
import numpy as np
from torch_geometric.datasets import Planetoid
from torch_geometric.transforms import NormalizeFeatures

from models.gnn import GraphEncoder
from retrieval.index import FaissIndex
from retrieval.search import GraphRetriever
from utils.metrics import precision_at_k, recall_at_k, mrr

# ============================================================
# Step 1: Load Dataset
# ============================================================
print("Loading Cora dataset...")
dataset = Planetoid(root='data', name='Cora', transform=NormalizeFeatures())
data = dataset[0]

print(f"Number of nodes: {data.num_nodes}")
print(f"Number of edges: {data.num_edges}")
print(f"Node feature dimension: {data.num_node_features}")
print(f"Number of classes: {dataset.num_classes}")

# ============================================================
# Step 2: Train GNN Encoder
# ============================================================
print("\nTraining GNN encoder...")

device = 'cuda' if torch.cuda.is_available() else 'cpu'
print(f"Using device: {device}")

# Initialize model
model = GraphEncoder(
    input_dim=data.num_node_features,
    hidden_dim=128,
    output_dim=64,
    num_layers=2,
    gnn_type='gcn'
)
model = model.to(device)

# Move data to device
x = data.x.to(device)
edge_index = data.edge_index.to(device)
y = data.y.to(device)

# Simple training: use node labels for supervision
# (In practice, you'd use contrastive learning)
optimizer = torch.optim.Adam(model.parameters(), lr=0.01)
criterion = torch.nn.CrossEntropyLoss()

# Add a classification head for training
classifier = torch.nn.Linear(64, dataset.num_classes).to(device)

model.train()
for epoch in range(200):
    optimizer.zero_grad()

    embeddings = model(x, edge_index)
    logits = classifier(embeddings)

    # Use only training nodes
    loss = criterion(logits[data.train_mask], y[data.train_mask])
    loss.backward()
    optimizer.step()

    if (epoch + 1) % 50 == 0:
        model.eval()
        with torch.no_grad():
            embeddings = model(x, edge_index)
            logits = classifier(embeddings)
            pred = logits.argmax(dim=1)

            train_acc = (pred[data.train_mask] == y[data.train_mask]).float().mean()
            test_acc = (pred[data.test_mask] == y[data.test_mask]).float().mean()

        print(f"Epoch {epoch+1}: Loss={loss:.4f}, Train Acc={train_acc:.4f}, Test Acc={test_acc:.4f}")
        model.train()

# ============================================================
# Step 3: Generate Embeddings for All Nodes
# ============================================================
print("\nGenerating embeddings...")

model.eval()
with torch.no_grad():
    all_embeddings = model(x, edge_index).cpu().numpy()

print(f"Embedding shape: {all_embeddings.shape}")

# ============================================================
# Step 4: Build FAISS Index
# ============================================================
print("\nBuilding FAISS index...")

index = FaissIndex(dimension=64, index_type='flat')
index.add(all_embeddings, ids=list(range(data.num_nodes)))

print(f"Index size: {index.index.ntotal}")

# ============================================================
# Step 5: Evaluate Retrieval
# ============================================================
print("\nEvaluating retrieval...")

# Ground truth: papers with same label are "relevant"
node_labels = data.y.cpu().numpy()

def get_relevant_docs(query_id):
    """Get all nodes with same label as query."""
    query_label = node_labels[query_id]
    return set(np.where(node_labels == query_label)[0]) - {query_id}

# Evaluate on test nodes
test_nodes = torch.where(data.test_mask)[0].cpu().numpy()
all_retrieved = []
all_relevant = []

for query_id in test_nodes[:100]:  # Evaluate 100 queries
    query_emb = all_embeddings[query_id].reshape(1, -1)
    scores, ids = index.search(query_emb, k=20)

    # Remove self from results
    retrieved = [int(i) for i in ids[0] if i != query_id][:10]
    relevant = get_relevant_docs(query_id)

    all_retrieved.append(retrieved)
    all_relevant.append(relevant)

# Calculate metrics
p_at_5 = np.mean([precision_at_k(r, rel, 5) for r, rel in zip(all_retrieved, all_relevant)])
p_at_10 = np.mean([precision_at_k(r, rel, 10) for r, rel in zip(all_retrieved, all_relevant)])
r_at_10 = np.mean([recall_at_k(r, rel, 10) for r, rel in zip(all_retrieved, all_relevant)])
mrr_score = mrr(all_retrieved, all_relevant)

print(f"\nRetrieval Results:")
print(f"  Precision@5:  {p_at_5:.4f}")
print(f"  Precision@10: {p_at_10:.4f}")
print(f"  Recall@10:    {r_at_10:.4f}")
print(f"  MRR:          {mrr_score:.4f}")

# ============================================================
# Step 6: Interactive Demo
# ============================================================
print("\n" + "="*50)
print("Interactive Demo")
print("="*50)

# Pick a random query
query_id = np.random.choice(test_nodes)
query_emb = all_embeddings[query_id].reshape(1, -1)
scores, ids = index.search(query_emb, k=6)

print(f"\nQuery: Node {query_id} (Label: {node_labels[query_id]})")
print(f"\nTop 5 similar nodes:")
for i, (score, node_id) in enumerate(zip(scores[0][1:6], ids[0][1:6])):  # Skip self
    match = "MATCH" if node_labels[node_id] == node_labels[query_id] else "     "
    print(f"  {i+1}. Node {node_id:4d} (Label: {node_labels[node_id]}) - Score: {score:.4f} {match}")

print("\nDone!")
\end{lstlisting}

\begin{explanation}
\textbf{Lines 20-27:} Load the Cora citation network. Each node is a paper with bag-of-words features. Edges are citations.

\textbf{Lines 38-46:} Initialize GNN with 2 layers, 128 hidden dim, 64 output dim.

\textbf{Lines 54-60:} For this example, we train with node classification as supervision. The embeddings learn to represent paper topics.

\textbf{Lines 88-92:} Generate embeddings for all nodes in one forward pass.

\textbf{Lines 98-100:} Build FAISS index with all document embeddings.

\textbf{Lines 107-109:} Define ground truth: papers with the same label (topic) are considered relevant.

\textbf{Lines 112-123:} Evaluate: for each query, retrieve top-10, compare to ground truth.

\textbf{Lines 140-148:} Interactive demo showing actual retrieval results.
\end{explanation}

\begin{checkpoint}
Run the complete example:
\begin{lstlisting}[language=bash]
python example.py
\end{lstlisting}

You should see:
\begin{itemize}
    \item[$\square$] Dataset loads (2708 nodes, 10556 edges)
    \item[$\square$] Training progresses (loss decreases)
    \item[$\square$] Test accuracy > 0.70
    \item[$\square$] Retrieval metrics printed
    \item[$\square$] Demo shows similar papers
\end{itemize}
\end{checkpoint}

\begin{gitcommit}
\begin{lstlisting}[language=bash]
git add .
git commit -m "Add complete end-to-end example on Cora dataset"
git push
\end{lstlisting}
\end{gitcommit}

%==============================================================================
% APPENDIX
%==============================================================================
\appendix

\section{Mathematical Notation Reference}

\begin{tabular}{|l|l|}
\hline
\textbf{Symbol} & \textbf{Meaning} \\
\hline
$G = (V, E)$ & Graph with vertices V and edges E \\
$A$ & Adjacency matrix \\
$D$ & Degree matrix \\
$\mathbf{h}_v$ & Embedding of node $v$ \\
$\mathcal{N}(v)$ & Neighbors of node $v$ \\
$\mathbf{W}$ & Learnable weight matrix \\
$\sigma$ & Activation function \\
$||\cdot||$ & Vector norm (magnitude) \\
$\mathbf{a} \cdot \mathbf{b}$ & Dot product \\
\hline
\end{tabular}

\section{Code Reference}

All files created:
\begin{itemize}
    \item \texttt{models/gnn.py} - GNN encoder models
    \item \texttt{train.py} - Training loop
    \item \texttt{retrieval/index.py} - FAISS index
    \item \texttt{retrieval/search.py} - Retrieval pipeline
    \item \texttt{utils/metrics.py} - Evaluation metrics
    \item \texttt{example.py} - Complete example
\end{itemize}

\section{Further Reading}

\begin{enumerate}
    \item \textbf{DeepWalk}: Perozzi et al., ``DeepWalk: Online Learning of Social Representations'' (KDD 2014)
    \item \textbf{Node2Vec}: Grover \& Leskovec, ``node2vec: Scalable Feature Learning for Networks'' (KDD 2016)
    \item \textbf{GCN}: Kipf \& Welling, ``Semi-Supervised Classification with Graph Convolutional Networks'' (ICLR 2017)
    \item \textbf{GraphSAGE}: Hamilton et al., ``Inductive Representation Learning on Large Graphs'' (NeurIPS 2017)
    \item \textbf{GAT}: Veli\v{c}kovi\'c et al., ``Graph Attention Networks'' (ICLR 2018)
\end{enumerate}

\vspace{2cm}
\begin{center}
\textit{Congratulations on building a graph-based retrieval system!}
\end{center}

\end{document}
